\begin{Theorem}
\label{thm:hilbert-basis-equivalent-characterizations}
Let $\{x_i\}_{i\in I}$ be an orthonormal set in a Hilbert space $H$. Then the following are equivalent:\begin{enumerate}[label=(\roman*)]
\item $\{x_i\}_{i\in I}$ is a basis.
\item For every $x,y\in H$, we have
\begin{equation}\label{eq:hilbert-basis-inner-product-expansion}
(x,y)=\sum_{i\in I}(x,x_i)(x_i,y).
\end{equation}

\item For every $x\in X$, we have
\begin{equation}\label{eq:hilbert-basis-parseval-characterization}
\|x\|^2=\sum_{i\in I}|(x_i,x)|^2.
\end{equation}

\item The set $\{x_i\}_{i\in I}$ is maximal, i.e., there is no orthonormal set containing it as a proper subset. In other words, if $x$ is orthogonal to each $x_i$, then $x=0$.\end{enumerate}
\end{Theorem}

\begin{Proof}
(i)$\Rightarrow$(ii): We have
\begin{equation}\label{eq:hilbert-basis-double-sum-inner-product}
\begin{aligned}
(x,y)&=\left(\sum_{i\in I}(x_i,x)x_i,\sum_{j\in I}(x_j,y)x_j\right)\\
&=\sum_{i,j\in I}(x,x_i)(x_j,y)\delta_{ij}\\
&=\sum_{i\in I}(x,x_i)(x_i,y).
\end{aligned}
\end{equation}
The exchange of summation and inner product is justified in the usual way by considering finite sums and then passing to the limit.

(ii)$\Rightarrow$(iii): Set $y=x$.

(iii)$\Rightarrow$(iv): If $x$ is orthogonal to each $x_i$, then \eqref{eq:hilbert-basis-parseval-characterization} implies $x=0$.

(iv)$\Rightarrow$(i): Let
\begin{equation}\label{eq:orthonormal-family-expansion-subspace}
Y:=\left\{x\in H\mid x=\sum_{i\in I}(x_i,x)x_i\right\}.
\end{equation}
Let $x^{(n)}$ be a Cauchy sequence in $Y$. Then there are at most countably many $i\in I$ for which $(x_i,x^{(n)})\ne0$ for any $n$. Let $\widetilde I$ be this at most countable set and let
\begin{equation}\label{eq:countable-orthonormal-expansion-subspace}
\widetilde Y:=\left\{x\in H\mid x=\sum_{i\in\widetilde I}(x_i,x)x_i\right\}.
\end{equation}
Parseval\textquotesingle s equality shows that $\widetilde Y$ is either finite-dimensional or isometric to the sequence space $\ell^2$ and hence complete; see Example~\ref{ex:ell-two-hilbert-space}. Therefore, the Cauchy sequence $x^{(n)}$ has a limit in $\widetilde Y\subseteq Y$, i.e., $Y$ is a closed subspace of $H$. On the other hand, (iv) says that $Y^\perp=\{0\}$, and by Theorem~\ref{thm:hilbert-orthogonal-projection} we conclude that $Y=H$.
\end{Proof}

\begin{Corollary}
\label{cor:hilbert-basis-existence}
Every Hilbert space has a basis.
\end{Corollary}

\begin{Proof}
A standard application of Zorn\textquotesingle s lemma shows that there is a maximal orthonormal set.
\end{Proof}

For separable Hilbert spaces, a basis can be found in a more constructive way using the Schmidt orthogonalization procedure. Let $\{x_n\}_{n\in\mathbb{N}}$ be a countable dense set. We then drop from this sequence each element which can be represented as a linear combination of the preceding ones. We thus end up with a new sequence $\{y_n\}$ of linearly independent elements such that the linear span of the $y_n$ is still dense in $H$. We now construct a sequence $z_n$ as follows:
\begin{equation}\label{eq:schmidt-orthogonalization-recursion}
\begin{aligned}
z_1&=y_1/\|y_1\|,\\
u_n&=y_n-\sum_{i=1}^{n-1}(z_i,y_n)z_i,\\
z_n&=u_n/\|u_n\|.
\end{aligned}
\end{equation}
It is easy to see that the $z_n$ are orthonormal, and their linear span is the same as that of the $y_n$, hence dense in $H$. Hence (iv) of Theorem~\ref{thm:hilbert-basis-equivalent-characterizations} applies and the $z_n$ form a basis.
